\chapter{Clustering and Dimension Reduction}

\section{Unsupervised methods}

Unsupervised learning finds structure in data without a target variable.
\hayashi{} provides clustering algorithms, non-linear embeddings, and
PCA-based tools.

\section{K-Means}

K-Means partitions observations into $K$ spherical clusters. It is fast
and scales well, but it assumes clusters of similar size and shape.

\begin{lstlisting}[caption={K-Means}]
set_seed(42)
let n = 300
let c1 = [0, 0]
let c2 = [5, 0]
let c3 = [2.5, 4]
generate df x1 = rnormal(n, 0, 0.7)
generate df x2 = rnormal(n, 0, 0.7)

let m_km = kmeans(df, x="x1,x2", k=3)
print(m_km)
\end{lstlisting}

Use `m_km.labels` to read the assigned cluster and `m_km.centroids` to
inspect the cluster means.

\section{Hierarchical clustering and DBSCAN}

Hierarchical clustering builds a tree of nested clusters. `hclust`
can be cut at any number of groups.

\begin{lstlisting}[caption={Hierarchical clustering}]
let m_hc = hclust(df, x="x1,x2", k=3, linkage="ward")
print(m_hc)
\end{lstlisting}

DBSCAN finds dense regions and labels sparse observations as noise.

\begin{lstlisting}[caption={DBSCAN}]
let m_db = dbscan(df, x="x1,x2", eps=0.8, min=5)
print(m_db)
\end{lstlisting}

\section{Gaussian Mixture Models}

GMM clustering models each cluster as a multivariate Gaussian. It is a
soft variant of K-Means that allows elliptical clusters.

\begin{lstlisting}[caption={GMM clustering}]
let m_gmm = gmm_clust(df, x="x1,x2", k=3)
print(m_gmm)
\end{lstlisting}

\section{Non-linear embeddings}

t-SNE and UMAP reduce high-dimensional data to two or three dimensions
while preserving local structure. They are useful for visualisation.

\begin{lstlisting}[caption={t-SNE}]
let m_tsne = tsne(df, x="x1,x2,x3", k=2, perplexity=30)
print(m_tsne)
\end{lstlisting}

\begin{lstlisting}[caption={UMAP}]
let m_umap = umap(df, x="x1,x2,x3", k=2, neighbors=15)
print(m_umap)
\end{lstlisting}

\section{Spectral clustering}

Spectral clustering uses the eigenstructure of a similarity graph. It
can capture non-convex cluster shapes where K-Means fails.

\begin{lstlisting}[caption={Spectral clustering}]
let m_sc = spectral(df, x="x1,x2", k=3)
print(m_sc)
\end{lstlisting}

\section{PCA, factor, and biplot}

PCA and factor analysis reduce dimensionality linearly. `biplot`
visualises observations and variables in the same PCA space.

\begin{lstlisting}[caption={PCA}]
let m_pca = pca(df, x="x1,x2,x3", n=2)
print(m_pca)
\end{lstlisting}

\begin{lstlisting}[caption={Biplot}]
let m_biplot = biplot(df, x="x1,x2,x3", type="symmetric")
print(m_biplot)
\end{lstlisting}

\section{Kernel density estimation}

`kde` estimates the density of a single variable with a chosen kernel.

\begin{lstlisting}[caption={KDE}]
let m_kde = kde(df, x1, bw=0.5, kernel="gaussian")
print(m_kde)
\end{lstlisting}
